% ================= APPENDIX: CAD & CONSTRUCTION =================
\section{CAD and Construction Detail}
\label{app:cad}

This appendix holds the mechanical detail deliberately kept out of
Section~\ref{sec:mechanical}: the model tree, the drawing set, the fabrication
settings and the assembly sequence. The rationale for the architecture --- why
the frame is split, what fixes the envelope, how the wheel is tuned --- stays in
Section~\ref{sec:mechanical} and is not repeated here. Anything in this appendix
may change without the design intent changing; anything in
Section~\ref{sec:mechanical} may not.

\subsection{Design Parameter Table}
\label{app:cad_params}

Table~\ref{tab:cad_params} is the single source of the driving dimensions
referenced by every part in the model, per convention~1 of
Section~\ref{sec:cad_method}. It is the CAD-side mirror of the envelope closure
of Section~\ref{sec:massbudget} and is updated on each iteration.

\begin{table}[h]
\centering
\caption{Driving CAD parameters. Every dependent feature references these; no
part carries a duplicated hard-coded value.}
\label{tab:cad_params}
\small
\begin{tabular}{l L{6.2cm} c}
\toprule
Symbol & Parameter & Value \\
\midrule
$L$         & Cube edge length (outer)                & \qty{180}{\milli\metre} \\
$D_w$       & Wheel outer diameter (Sec.~\ref{sec:rw_sizing}) & \qty{130}{\milli\metre} \\
$w_r$       & Wheel ring radial width                 & \qty{20}{\milli\metre} \\
$t_r$       & Wheel ring axial thickness              & \qty{12}{\milli\metre} \\
$r_b$       & Ballast pocket circle radius            & \qty{55}{\milli\metre} \\
--          & Board outline (Sec.~\ref{sec:layout})   & \qtyproduct{150 x 90}{\milli\metre} \\
$t_w$       & Frame wall thickness                    & \textcolor{red}{TBD} \\
$d_m$       & Motor axis offset from cube centre      & \textcolor{red}{TBD} \\
$h_m$       & Motor mounting face height on its axis  & \textcolor{red}{TBD} \\
$c_w$       & Wheel-to-subframe axial clearance       & \textcolor{red}{TBD} \\
$c_o$       & Clearance between orthogonal wheel modules & \textcolor{red}{TBD} \\
$V_{\min}$  & Minimum internal clear volume           & \textcolor{red}{TBD} \\
$r_c$       & Corner contact feature radius           & \textcolor{red}{TBD} \\
\bottomrule
\end{tabular}
\end{table}

\paragraph{Open input.} The frame wall thickness and internal clearance that
set the ``$L$ minus \qtyrange{30}{40}{\milli\metre}'' wheel-diameter allowance
used in Table~\ref{tab:envelope_trade} are an estimate rather than a CAD result.
The envelope argument holds anywhere in that band, but the selected ring
diameter depends on where in it the real frame lands, so the maximum wheel
diameter that actually fits the \qty{180}{\milli\metre} frame must be confirmed
from the model before the ring is frozen.

\subsection{Model Tree and Part List}
\label{app:cad_tree}

Table~\ref{tab:cad_parts} lists the modelled parts by assembly, with their
fabrication route and current status. It is the mechanical counterpart to the
electronics bill of materials of Appendix~\ref{app:bom}: the parts here are the
team's design responsibility and are not supplied.

\begin{table}[h]
\centering
\caption{CAD model tree and part list. Quantities are per assembled cube unless
noted; the 1-DoF rig parts are listed separately.}
\label{tab:cad_parts}
\small
\setlength{\tabcolsep}{4pt}
\begin{tabular}{c L{4.0cm} c L{3.4cm} L{3.0cm}}
\toprule
Ref & Part & Qty & Fabrication route & Status \\
\midrule
\multicolumn{5}{@{}l}{\textbf{Inner structure}}\\
\midrule
M1 & Motor subframe             & \textcolor{red}{TBD} & \textcolor{red}{TBD} & \textcolor{red}{TBD} \\
M2 & Motor mount ($\times$axis) & 3 & \textcolor{red}{TBD} & \textcolor{red}{TBD} \\
M3 & Battery cradle / retention & 1 & \textcolor{red}{TBD} & \textcolor{red}{TBD} \\
M4 & Electronics tray           & 1 & \textcolor{red}{TBD} & \textcolor{red}{TBD} \\
M5 & Brake servo bracket        & 3 & \textcolor{red}{TBD} & \textcolor{red}{TBD} \\
\midrule
\multicolumn{5}{@{}l}{\textbf{Reaction wheel assembly}}\\
\midrule
W1 & Ballasted ring, \qty{130}{\milli\metre} OD & 3 & Printed PET-CF & Geometry fixed \\
W2 & Wheel hub / rotor adapter  & 3 & Printed & \textcolor{red}{TBD} \\
W3 & M6 hex nut ballast (ISO 4032) & 45 & Standard hardware & 15 per wheel \\
W4 & Ballast retaining cover    & 3 & \textcolor{red}{TBD} & Required, not yet designed \\
W5 & Brake strike feature       & 3 & \textcolor{red}{TBD} & \textcolor{red}{TBD} \\
\midrule
\multicolumn{5}{@{}l}{\textbf{Outer frame}}\\
\midrule
F1 & Frame member / edge rail   & \textcolor{red}{TBD} & \textcolor{red}{TBD} & \textcolor{red}{TBD} \\
F2 & Corner contact feature     & 8 & \textcolor{red}{TBD} & \textcolor{red}{TBD} \\
F3 & Face panel / wheel guard   & \textcolor{red}{TBD} & \textcolor{red}{TBD} & \textcolor{red}{TBD} \\
\midrule
\multicolumn{5}{@{}l}{\textbf{1-DoF rig (Sec.~\ref{sec:cad_2d})}}\\
\midrule
R1 & Face plate                 & 1 & \textcolor{red}{TBD} & \textcolor{red}{TBD} \\
R2 & Pivot bearing mount        & 1 & \textcolor{red}{TBD} & \textcolor{red}{TBD} \\
R3 & Rig motor mount            & 1 & \textcolor{red}{TBD} & \textcolor{red}{TBD} \\
R4 & Hard stops / containment   & \textcolor{red}{TBD} & \textcolor{red}{TBD} & \textcolor{red}{TBD} \\
\bottomrule
\end{tabular}
\end{table}

\subsection{2D Prototype Construction}
\label{app:cad_2d}

% Rig drawings and construction notes: plate geometry, pivot bearing seat,
% motor mount, brake barrier travel, hard-stop placement, and the mounting
% location of the IMU relative to the pivot (Sec.~\ref{sec:control2d}).

% Figure placeholder: rig general-assembly view with the pivot, motor axis and
% brake labelled.

\subsection{3D Cube Construction}
\label{app:cad_3d}

% Cube drawings: inner structure general assembly, the three motor modules on
% their orthogonal axes, outer frame and contact features, and the section view
% showing the battery at the geometric centre.

% Figure placeholder: cube exploded view.
% Figure placeholder: section view through two orthogonal wheel modules,
% showing the clearances c_w and c_o of Table~\ref{tab:cad_params}.

\paragraph{Wheel pocket placement --- a CAD constraint the sizing cannot
check.} The ballast pockets must be drawn in three groups of five within the
inter-spoke sectors, \emph{not} at a uniform angular pitch. The reasoning and
the angular figures are in Section~\ref{sec:rw_sizing}; what matters here is
that the sizing calculation models the ring as a bare annulus with no knowledge
of the spokes, so it will report a feasible design for a uniform layout that in
fact places pockets on the spoke roots. The layout rule is
Table~\ref{tab:pocket_layout} and it governs the drawing.

\begin{table}[h]
\centering
\caption{Ballast pocket layout rule. Grouping by sector preserves three-fold
symmetry and caps the ring at 21 pockets.}
\label{tab:pocket_layout}
\small
\begin{tabular}{L{7.0cm} c}
\toprule
Quantity & Value \\
\midrule
Spoke half-angle at ring inner diameter & \ang{6.38} \\
Pocket half-angle at $r_b = \qty{55}{\milli\metre}$ & \ang{6.27} \\
Required spoke-centre to pocket-centre separation & \ang{12.65} \\
Usable arc per sector, after spoke keep-outs & \ang{94.7} \\
Pockets per sector at $N = 15$ & 5 \\
Maximum pockets per sector & 7 (21 total) \\
\bottomrule
\end{tabular}
\end{table}

\subsection{Fabrication and Assembly}
\label{app:cad_fab}

\subsubsection{Print and Manufacturing Settings}

The printed parts are structural, not cosmetic: the wheel ring carries
centrifugal load at the ballast pockets and the motor mounts carry the brake
impulse, which is the design load case identified in
Section~\ref{sec:background}. Settings are therefore recorded per part in
Table~\ref{tab:print_settings} and treated as part of the design, since a part
reprinted at a different orientation or infill is not the part that was
analysed.

\begin{table}[h]
\centering
\caption{Fabrication settings for the structural printed parts. Layer
orientation is specified relative to the principal load direction.}
\label{tab:print_settings}
\small
\begin{tabular}{l L{2.6cm} L{2.8cm} L{4.0cm}}
\toprule
Ref & Material & Infill / walls & Orientation and rationale \\
\midrule
W1 & PET-CF, \qty{1290}{\kilo\gram\per\metre\cubed} & \textcolor{red}{TBD} & Layers normal to the spin axis, so the hoop load and the ballast pocket walls are not carried across a layer boundary \\
W2 & \textcolor{red}{TBD} & \textcolor{red}{TBD} & Carries the brake impulse into the rotor; orientation set by that load path \\
M1 & \textcolor{red}{TBD} & \textcolor{red}{TBD} & \textcolor{red}{TBD} \\
M2 & \textcolor{red}{TBD} & \textcolor{red}{TBD} & Carries the brake reaction; the governing load case of Sec.~\ref{sec:background} \\
F1 & \textcolor{red}{TBD} & \textcolor{red}{TBD} & \textcolor{red}{TBD} \\
F2 & \textcolor{red}{TBD} & \textcolor{red}{TBD} & \textcolor{red}{TBD} \\
\bottomrule
\end{tabular}
\end{table}

The wheel ring density is not a free choice at print time: the inertia figures
of Section~\ref{sec:rw_sizing} assume
\qty{1290}{\kilo\gram\per\metre\cubed} throughout, so a ring printed at reduced
infill supplies less inertia than designed and moves the nut count. Print
settings for W1 are therefore part of the sizing, and a re-print at different
settings requires the ballast count to be re-derived.

\subsubsection{Assembly Sequence}

The sequence of Table~\ref{tab:assembly_seq} is ordered by access: each step is
one whose verification becomes impossible or destructive once a later step is
complete. The encoder air gap in particular is re-checked after mechanical
assembly, since it can shift when the structure is bolted up
(Table~\ref{tab:signal_rules}).

\begin{table}[h]
\centering
\caption{Assembly sequence. Ordering is set by access: each check must be
completed before the step that encloses it.}
\label{tab:assembly_seq}
\small
\begin{tabular}{c L{5.6cm} L{6.4cm}}
\toprule
Step & Operation & Check before proceeding \\
\midrule
1 & Wheel assembly: ring, hub, ballast bolts to the selected count & Balance check; bolt retention per Sec.~\ref{sec:rw_sizing} \\
2 & Motor to mount, encoder magnet to rotor & Encoder air gap and centring to the assembly spec \\
3 & Motor modules to inner structure, three axes & Axis orthogonality; wheel clearance $c_w$, $c_o$ \\
4 & Battery and electronics tray into inner structure & Pack centring; board keep-out clearances \\
5 & Harness install and dressing & Continuity and isolation (Sec.~\ref{sec:electrical_safety}) \\
6 & Inner structure into outer frame & Encoder air gaps re-checked after bolt-up \\
7 & Contact features and guards & Contact geometry; wheel guard clearance \\
8 & Mass properties measurement & Measured mass against Table~\ref{tab:massbudget} \\
\bottomrule
\end{tabular}
\end{table}

\subsection{Mass Properties Verification}
\label{app:cad_mass}

The CAD roll-up is an estimate until the assembly is weighed. Step~8 above
closes the loop of Section~\ref{sec:massbudget}: the measured mass is compared
against the modelled value, and any discrepancy is carried back into
$I_{w,\text{target}}$ and absorbed, where possible, by the ballast trim of
Section~\ref{sec:rw_sizing} rather than by a re-design.

\begin{table}[h]
\centering
\caption{Modelled versus measured mass properties. A discrepancy here is
absorbed by ballast trim where the margin allows.}
\label{tab:mass_verification}
\small
\begin{tabular}{L{4.6cm} c c c}
\toprule
Quantity & Design value & Measured & Discrepancy \\
\midrule
Total cube mass $M$          & \qty{1400}{\gram} & \textcolor{red}{TBD} & \textcolor{red}{TBD} \\
Wheel mass $m_w$ (each)      & \qty{156}{\gram}  & \textcolor{red}{TBD} & \textcolor{red}{TBD} \\
Wheel inertia $I_w$ (each)   & $4.341\times10^{-4}$ & \textcolor{red}{TBD} & \textcolor{red}{TBD} \\
Ballast nut mass (each)      & \qty{2.5}{\gram}  & \textcolor{red}{TBD} & \textcolor{red}{TBD} \\
CoM offset from cube centre  & 0 (by design)     & \textcolor{red}{TBD} & \textcolor{red}{TBD} \\
\bottomrule
\end{tabular}
\end{table}

Two of these rows are gating rather than informational. The measured $M$ decides
whether the mass budget of Section~\ref{sec:massbudget} closes at all, and it
propagates through $h_w \propto M$ into the wheel; the measured nut mass
propagates into $I_w$ directly. On the latter, the design uses a measured
\qty{2.5}{\gram} against a geometric estimate of \qty{2.617}{\gram} for an M6
nut to ISO 4032 --- a \qty{-4.5}{\percent} discrepancy consistent with ordinary
manufacturing tolerance and plating variation, but worth re-confirming on a
scale before the nut count is frozen, since the shape (and hence the radius of
gyration) is taken from the nominal geometry even when the mass is measured.
